% ---- pg-461 （书页 449） ----

又使用两种说法: 将观测值指派给各总体, 或者将对象指派给各类别.

下面回顾一下第 1 章中所介绍的一些分类问题的例子.

\begin{mvtable}
\begin{tabular}{@{}p{0.30\linewidth}p{0.62\linewidth}@{}}
\toprule
总体 $\pi_1$ 和 $\pi_2$ & 所测量的变量 $\mathbf{X}$\\
\midrule
1. 有偿付力与无偿付力的财产责任保险公司. & 总资产, 股票与债券价值, 股票与债券的市值, 损失支出, 盈余, 签订的保费金额.\\
2. 非溃疡胃病组(胃功能紊乱者)与控制组(“正常”者). & 焦虑、依赖性、罪恶感、完美主义的量度.\\
3. 由麦迪逊和汉密尔顿所撰写的联邦文献. & 不同单词的出现频数, 句子的长度.\\
4. 两种鸢尾(译者注: 野花). & 萼片与花瓣的长度, 花瓣裂缝的深度, 苞的长度, 花粉直径.\\
5. 新产品的速购者与迟购者. & 教育, 收入, 家庭大小, 过去更换品牌的次数.\\
6. 成功与不成功(肄业)的大学生. & 入学考试成绩, 高中毕业平均成绩, 高中活动的数目.\\
7. 男性与女性. & 人类学指标测量值, 相似环境及古颅骨体积.\\
8. 良好信用与不良信用风险. & 收入, 年龄, 信用卡数目, 家庭规模.\\
9. 嗜酒者与非嗜酒者. & 单胺氧化酶的活性, 环脘酶的活性.\\
\bottomrule
\end{tabular}
\end{mvtable}

例如, 我们从第 5 项中看到, 对象(消费者)要被分成两类(“速购者”与“迟购者”), 分类的依据是设想的有关变量的观测值(教育、收入等等). 用观测值和总体这样的术语, 我们想识别形如 $\mathbf{x}'=[x_1(\text{教育}),x_2(\text{收入}),x_3(\text{家庭规模}),x_4(\text{更换品牌次数})]$ 的观察值是属于总体 $\pi_1$(速购者), 还是属于 $\pi_2$(迟购者).

目前, 我们集中讨论两个总体的分类, 到 11.3 节中再回来讨论分离.

分配或分类法则通常由“学习”样本发展而来. 对已知随机选自两个总体的对象测出特征, 并考察这些特征之间有无差别. 从本质上说, 全部可能的样本结果的集合被分成 $R_1$ 和 $R_2$ 两个区域. 要是某个新观测值落入 $R_1$, 就将它分配给总体 $\pi_1$; 若落入 $R_2$, 就将它分配给总体 $\pi_2$. 这样, 一个观测值集合倾向于 $\pi_1$, 而另一个集合则倾向于 $\pi_2$.

在这一点上也许会感到奇怪: 我们怎么知道某些观测值属于某个特定总体而对其余观测值却不能肯定? (当然, 正是这一点才使分类成为一个需要研究的课题!)有几种情况会造成这种表现异常(见文献[20]):

\noindent\mvbold{1. 关于未来表现的不完备的知识}

例: 在一个企业破产之前的两年内, 曾对某些财务变量的极端值进行了观测. 根据这些领先指标的观测值, 另一个企业的官员便可对本企业的状态是否健康作出判断, 并在必要时采取及时的纠正行动.

医学院的研究生招生办公室可以根据报名学生的考试成绩和其他记录把报名学生分成“有可能取得硕士学位”和“不大可能取得硕士学位”两类. 这里, 真正的结果只能在几年的研究生培训结束之后才能确定.

\noindent\mvbold{2. “完备的”信息需要毁坏对象}

例: 计算器电池的寿命只能通过使用直到将它用坏时才能确定; 一块木材的强度只能通过在其上加荷, 直到它断裂时才能得到; 已被毁坏的产品不能出售. 为避免上述情况, 人们希望根据产品的某些事先测量值将它们分成“好的”或“坏的”.

\noindent\mvbold{3. 不可获得的或获得成本昂贵的信息}

例: 某些联邦文件可断定是由詹姆斯·麦迪逊或亚历山大·汉密尔顿撰写的, 因为这些
